\subsection{Fundamentação Técnica dos Conceitos Utilizados}
\label{sec:fundamentacao}

Antes de apresentar as tecnologias específicas adotadas, esta subseção
fundamenta os principais conceitos de Engenharia de Software que orientaram
as decisões de projeto do Movy. O objetivo é situar o leitor quanto aos
referenciais teóricos e normativos que sustentam as escolhas descritas nas
subseções seguintes, evitando que a discussão técnica se reduza a uma
enumeração de ferramentas.

\subsubsection{Requisitos de Software}

Requisitos de software descrevem as funções, restrições e qualidades
esperadas de um sistema, servindo de contrato entre as necessidades dos
interessados e a solução a ser construída. A literatura clássica distingue
\textbf{requisitos funcionais} --- que especificam o que o sistema deve fazer
--- de \textbf{requisitos não funcionais} --- que expressam atributos de
qualidade, como desempenho, segurança e usabilidade \cite{sommerville2011,
pressman2016}. A norma \textit{ISO/IEC/IEEE 29148} consolida as boas práticas
de engenharia de requisitos, recomendando que cada requisito seja verificável,
não ambíguo e rastreável \cite{ieee29148}. No presente trabalho, os requisitos
são tratados como \textbf{requisitos do sistema}, expostos pela API e
materializados em telas pela aplicação cliente, conforme detalhado na
Seção~\ref{sec:arquitetura}.

\subsubsection{Modelagem de Software}

A modelagem de software emprega abstrações gráficas para representar a
estrutura e o comportamento de um sistema antes e durante a sua construção.
A \textit{Unified Modeling Language} (UML), mantida pela \textit{Object
Management Group}, é a notação de referência para essa finalidade, oferecendo
diagramas como o de \textbf{casos de uso} (que captura as interações entre
atores e o sistema), o de \textbf{classes} (que descreve entidades e seus
relacionamentos) e modelos de dados \cite{omguml}. A modelagem antecipa
decisões de projeto, favorece a comunicação entre os envolvidos e reduz o
retrabalho na implementação \cite{sommerville2011}. Neste trabalho, a
modelagem abrange casos de uso, modelo de domínio (diagrama de classes) e o
modelo entidade-relacionamento (Seção~\ref{sec:arquitetura} e seguintes).

\subsubsection{Arquitetura em Camadas, Clean Architecture e DDD}

A organização de um sistema em \textbf{camadas} separa responsabilidades
distintas --- apresentação, aplicação, domínio e infraestrutura ---, de modo
que cada camada dependa apenas das mais internas, e não o contrário. A
\textit{Clean Architecture} formaliza esse princípio por meio da
\textbf{regra de dependência}: as regras de negócio centrais não conhecem
detalhes de \textit{frameworks}, banco de dados ou interface, tornando-se
testáveis e resistentes a mudanças tecnológicas \cite{martin2017}. O
\textit{Domain-Driven Design} (DDD) complementa essa abordagem ao colocar o
\textbf{modelo de domínio} no centro do projeto, propondo conceitos como
entidades, \textit{value objects}, agregados e contextos delimitados
(\textit{bounded contexts}) para alinhar o código à linguagem do negócio
\cite{evans2003}. Esses dois referenciais orientam tanto a estrutura do
servidor quanto, por espelhamento, a da aplicação cliente.

\subsubsection{APIs REST}

\textit{Representational State Transfer} (REST) é um estilo arquitetural para
sistemas distribuídos, definido por Fielding, que estabelece restrições como
a comunicação \textbf{cliente-servidor}, a ausência de estado de sessão no
servidor (\textit{statelessness}), a possibilidade de \textit{cache} e uma
interface uniforme baseada em recursos identificados por URIs e manipulados
por métodos HTTP \cite{fielding2000}. Uma API REST bem desenhada favorece a
escalabilidade horizontal, o desacoplamento entre cliente e servidor e a
integração com múltiplos consumidores. No Movy, a API REST concentra toda a
lógica de domínio e é consumida por uma aplicação cliente web, mas permanece
apta a atender outros clientes (por exemplo, aplicações móveis).

\subsubsection{SaaS e Multi-tenancy}

\textit{Software as a Service} (SaaS) é um modelo de entrega de software no
qual a aplicação é executada em infraestrutura do provedor e disponibilizada
aos clientes pela rede, sob demanda, sem instalação local \cite{nist800145}. É
comum que plataformas SaaS adotem a \textbf{multi-tenancy} (multitenância), em
que uma mesma instância da aplicação atende a diversas organizações
(\textit{tenants}) mantendo o \textbf{isolamento lógico} dos respectivos
dados. Esse isolamento pode ser obtido por bancos separados, esquemas
separados ou --- como neste trabalho --- por um discriminador de tenant em
cada registro (o campo \texttt{organizationId}), com filtragem aplicada em
todas as consultas e validada por mecanismos de autorização.

\subsubsection{Segurança, Autenticação e Autorização}

\textbf{Autenticação} é o processo de comprovar a identidade de um usuário,
ao passo que \textbf{autorização} determina o que esse usuário pode fazer.
Em APIs REST \textit{stateless}, é usual empregar \textit{tokens} assinados no
padrão \textit{JSON Web Token} (JWT), que transportam, de forma verificável,
as informações de identidade e permissões \cite{rfc7519}. As senhas devem ser
armazenadas apenas como \textit{hash} com função adaptativa e \textit{salt},
como o \textbf{bcrypt}, de modo a resistir a ataques de força bruta
\cite{bcrypt}. O \textbf{controle de acesso baseado em papéis} (RBAC) associa
permissões a papéis, e papéis a usuários, simplificando a administração das
autorizações \cite{nistrbac}. As diretrizes do \textit{OWASP} ---
particularmente o \textit{Top 10} e o \textit{Application Security
Verification Standard} (ASVS) --- orientam a prevenção de vulnerabilidades
recorrentes em aplicações web, como o \textit{Broken Access Control} e as
referências diretas inseguras a objetos (IDOR) \cite{owasptop10, owaspasvs}.
Por fim, o tratamento de dados pessoais observa os princípios da Lei Geral de
Proteção de Dados (LGPD) \cite{lgpd}.

\subsubsection{Testes de Software}

Testes de software constituem a atividade de verificação que executa o sistema
sob condições controladas para revelar defeitos e aumentar a confiança na sua
corretude \cite{pressman2016, sommerville2011}. Os \textbf{testes unitários}
exercitam, de forma isolada, as menores unidades de código (como casos de uso
e \textit{value objects}), favorecendo a detecção precoce de erros e
funcionando como documentação executável do comportamento esperado. Padrões
como o \textit{Arrange-Act-Assert} (AAA) e o uso de \textit{test doubles}
(\textit{mocks}) promovem testes legíveis e determinísticos. A estratégia de
verificação adotada neste trabalho é detalhada na Seção~\ref{subsec:verificacao}.
